%=============================================================================!
% 4.F  SOLUTIONS TO THE EXERCISES                                             !
%=============================================================================!
\unnumbered{Chapter~\ref{ch:oscillations}: Oscillations}
\label{sec:os-solutions}

\solutionto{ex:os-phase-pi}At the angle $\pi$ the point $P$ is at $(-A,
0)$, the left end of the circle's horizontal diameter, so the block is at
$x = A\cos\pi = -A$, a turning point, momentarily at rest. Going
anticlockwise, $P$ is moving straight down there, and next swings towards
the right, so the block is about to move back towards the centre, in the
positive direction.

\solutionto{ex:os-energy}$E = \frac12 kA^2 = \frac12\times80\times0.03^2 =
\SI{0.036}{\joule}$. With $\omega = \sqrt{80/0.2} =
\SI{20}{\radian\per\second}$, the largest speed is $A\omega = 0.03\times20 =
\SI{0.6}{\metre\per\second}$. At $x = \SI{0.015}{\metre}$, $U =
\frac12\times80\times0.015^2 = \SI{0.009}{\joule}$, so $K = 0.036 - 0.009
= \SI{0.027}{\joule}$ and
\begin{equation*}
   v = \sqrt{\frac{2\times0.027}{0.2}} = \SI{0.520}{\metre\per\second} .
\end{equation*}

\solutionto{ex:os-average}Rearranging $\cos2\theta = 1 - 2\sin^2\theta$
gives $\sin^2\theta = \frac12 - \frac12\cos2\theta$. With $\theta = \omega
t$, the average of $\cos 2\omega t$ over a period $2\pi/\omega$ is zero:
$\sin 2\omega t/(2\omega)$ has derivative $\cos 2\omega t$ by the chain
rule, so the integral is $[\sin 2\omega t/(2\omega)]_0^{2\pi/\omega} =
(\sin4\pi - \sin 0)/(2\omega) = 0$. The integral of a sum is the sum of
the integrals, and a constant averages to itself, so the average of
$\sin^2\omega t$ is $\frac12 - \frac12\times0 = \frac12$.

\solutionto{ex:os-ellipse}$E = \frac12 mv^2 + \frac12 kx^2 =
\frac12\times0.5\times0.4^2 + \frac12\times50\times0.03^2 = 0.04 + 0.0225 =
\SI{0.0625}{\joule}$. The semi-axes are $\sqrt{2E/k} = \sqrt{0.0025} =
\SI{5}{\centi\metre}$ and $\sqrt{2E/m} = \sqrt{0.25} =
\SI{0.5}{\metre\per\second}$. The velocity is negative, so $x$ is about to
shrink; the state lies in the lower half of the ellipse, moving left.

\solutionto{ex:os-clock}From $2\pi\sqrt{L/g} = 2$, squaring and solving
for $L$,
\begin{equation*}
   L = g\,\Bigl(\frac{2}{2\pi}\Bigr)^2 = \frac{9.80665}{\pi^2}
   = \SI{0.9936}{\metre} .
\end{equation*}

\solutionto{ex:os-arcsine}By the argument of \cref{eq:os-period}, the time
from $x = 0$ to $x = A/2$ is the integral of $\dd x/v$, with $v =
\omega\sqrt{A^2 - x^2}$ for the spring (\cref{sec:os-anharmonic}), so by
the toolbox on the derivative of an inverse function
\begin{equation*}
   t = \frac{1}{\omega}\int_0^{A/2}\frac{\dd x}{\sqrt{A^2 - x^2}}
   = \frac1\omega\Bigl[\arcsin\frac{x}{A}\Bigr]_0^{A/2}
   = \frac{1}{\omega}\Bigl(\arcsin\tfrac12 - 0\Bigr) = \frac{\pi}{6\omega} ,
\end{equation*}
since $\sin(\pi/6) = \frac12$: an equilateral triangle of side 1, cut in
half, has an angle of $\pi/6$ opposite a side of $\frac12$. For $\omega =
\SI{10}{\radian\per\second}$ this is $\pi/60 = \SI{0.0524}{\second}$, a
twelfth of the period $2\pi/\omega$. The whole way out, from the centre to
the turning point, takes a quarter period, $\pi/(2\omega) =
3\pi/(6\omega)$, so the second half of the way takes $3\pi/(6\omega) -
\pi/(6\omega) = 2\pi/(6\omega)$, twice as long as the first, because the
block slows down.

\solutionto{ex:os-double}With $s = x/a$, the solution of
\cref{ex:en-double} gives $U''(x) = (4U_0/a^2)(3s^2 - 1)$, which at $s =
\pm1$ is $8U_0/a^2$. By \cref{eq:os-small}, $\omega_0 = \sqrt{8U_0/(m
a^2)}$.

\solutionto{ex:os-isotropic}On the line $x = 0$, $\vb{r} = (0, y)$.
Writing $g$ for $\lvert\vb{g}\rvert$, the solution of \cref{ex:en-top}
gives $\vb{g}_1 = g(\frac{\sqrt3}{2}, -\frac12)$, $\vb{g}_2 = g(0, 1)$ and
$\vb{g}_3 = g(-\frac{\sqrt3}{2}, -\frac12)$, so the products
$\vb{g}_k\cdot\vb{r}$ are $-\frac12 gy$, $gy$ and $-\frac12 gy$. The
cosine is even, so
\begin{equation*}
      U(0, y) = \tfrac14 U_{\mathrm b}\bigl[3 - 2\cos(\tfrac12 gy)
   - \cos(gy)\bigr] .
\end{equation*}
Differentiating twice by the chain rule, $-2\cos(\frac12 gy)$ gives
$2(\frac12 g)^2\cos(\frac12 gy)$, which is $\frac12 g^2$ at $y = 0$, and
$-\cos(gy)$ gives $g^2\cos(gy)$, which is $g^2$. So $U''(0) =
\frac14U_{\mathrm b}(\frac12g^2 + g^2) = \frac38U_{\mathrm b}g^2$.
Squaring $g = 4\pi/(\sqrt3\,a)$ gives $g^2 = 16\pi^2/(3a^2)$, and
\begin{equation*}
   \tfrac38U_{\mathrm b}\,\frac{16\pi^2}{3a^2} = \frac{2\pi^2U_{\mathrm b}}{a^2} ,
\end{equation*}
the curvature along $x$: the well is equally stiff in both directions, and
by a longer argument in every direction.

\solutionto{ex:os-lj}By \cref{sec:os-anharmonic}, $U_{\mathrm{LJ}}'' =
(4\varepsilon/r^2)[156(\sigma/r)^{12} - 42(\sigma/r)^6]$, which is zero
when $156(\sigma/r)^6 = 42$, after dividing by $4\varepsilon/r^2$ and by
$(\sigma/r)^6$. So $(\sigma/r)^6 = 42/156 = 7/26$ and $r =
(26/7)^{1/6}\sigma = 1.2445\,\sigma$. The bracket is $(\sigma/r)^6\,
[156(\sigma/r)^6 - 42]$, positive at smaller $r$, where $(\sigma/r)^6$ is
larger, and negative at larger $r$, so inside this distance the slope of
$U$ steepens with $r$ and outside it flattens. Beyond $r_{\min}$ the force
$-U'$ pulls the atoms together with a size equal to that slope, so here
the attractive force is greatest.

\solutionto{ex:os-morse}Write $w = \mathrm{e}^{-a(r - r_0)}$, so that
$\dd w/\dd r = -aw$ by the chain rule, and $U_{\mathrm M} = D(1 - w)^2$.
Then
\begin{align*}
   \frac{\dd U_{\mathrm M}}{\dd r} &= 2D(1 - w)\Bigl(-\frac{\dd w}{\dd r}\Bigr)
   = 2Da\,w\,(1 - w) = 2Da\,(w - w^2) ,\\
   \frac{\dd^2 U_{\mathrm M}}{\dd r^2} &= 2Da\,(1 - 2w)\,\frac{\dd w}{\dd r}
   = -2Da^2\,w\,(1 - 2w) ,
\end{align*}
and at $r = r_0$, where $w = 1$, this is $-2Da^2\times1\times(-1) =
2Da^2$.

\solutionto{ex:os-liquid}$\omega_{\mathrm d} = \sqrt{10^2 - 6^2} =
\SI{8}{\radian\per\second}$, a period of $2\pi/8 = \SI{0.785}{\second}$.
The swing shrinks by $\mathrm{e}^{-\gamma t/2} = \mathrm{e}^{-6t}$, which
is $\frac1{10}$ when $6t = \ln10$, that is, taking the natural logarithm
of $\mathrm{e}^{6t} = 10$, at $t = 2.303/6 = \SI{0.384}{\second}$.

\solutionto{ex:os-critical}Write $y = c_1 + c_2t$ and $z =
\mathrm{e}^{-\omega_0 t}$ for short, so that $x = yz$, $\dot{y} = c_2$ and
$\dot{z} = -\omega_0 z$. By the product rule,
\begin{align*}
   \dot{x} &= z\,(c_2 - \omega_0 y) , \qquad
   \ddot{x} = z\,(-\omega_0 c_2) - \omega_0 z\,(c_2 - \omega_0 y)
   = z\,(\omega_0^2 y - 2\omega_0 c_2) .
\end{align*}
With $\gamma = 2\omega_0$, $\ddot{x} + \gamma\dot{x} + \omega_0^2 x =
z\,[\omega_0^2y - 2\omega_0c_2 + 2\omega_0c_2 - 2\omega_0^2y +
\omega_0^2y] = 0$, which is \cref{eq:os-damped} with its right side
moved to the left.

\solutionto{ex:os-off}By \cref{eq:os-resonance}, at $\omega = 5$,
\begin{equation*}
   A = \frac{1}{\sqrt{(100 - 25)^2 + (2\times5)^2}} = \frac{1}{75.66}\
   \si{\metre} = \SI{1.32}{\centi\metre} ,
      \qquad \tan\delta = \frac{10}{75} .
\end{equation*}
Here $\omega_0^2 - \omega^2 = 75$ is positive, so $\cos\delta$ and
$\sin\delta$ are both positive and $\delta$ lies between 0 and $\pi/2$,
inside the range of the arctangent: $\delta = \arctan(10/75) =
\SI{0.133}{\radian}$. At $\omega = 15$, $A = 1/\sqrt{125^2 + 30^2} =
1/128.55\ \si{\metre} = \SI{0.778}{\centi\metre}$ and $\tan\delta =
30/(-125)$. Now $\omega_0^2 - \omega^2 = -125$, so $\cos\delta$ is
negative while $\sin\delta$ is still positive, and $\delta$ lies between
$\pi/2$ and $\pi$, outside the range of the arctangent. By the rule of
\cref{sec:os-driven}, $\delta = \pi - \arctan(30/125) = \pi - 0.236 =
\SI{2.906}{\radian}$. Below resonance the block nearly follows the push;
above it, it nearly opposes it.

\solutionto{ex:os-linear}Let $x = c_1x_1 + c_2x_2$, and write
\cref{eq:os-damped} as $\ddot{x} + \gamma\dot{x} + \omega_0^2x = 0$,
adding $\omega_0^2x + \gamma\dot{x}$ to both sides. Derivatives of sums
are sums of derivatives, and constants pass through, so $\ddot{x} +
\gamma\dot{x} + \omega_0^2x = c_1(\ddot{x}_1 + \gamma\dot{x}_1 +
\omega_0^2x_1) + c_2(\ddot{x}_2 + \gamma\dot{x}_2 + \omega_0^2x_2) =
c_1\times0 + c_2\times0 = 0$, each bracket vanishing because $x_1$ and
$x_2$ are motions.

\solutionto{ex:os-isotope}The reduced masses are $m_{\mathrm{r},\mathrm{OH}} =
\SI{0.9483}{\amu}$ (\cref{sec:os-bodies}) and $m_{\mathrm{r},\mathrm{OD}} =
2.014\times15.999/18.013 = \SI{1.7888}{\amu}$. With the same stiffness,
$\omega = \sqrt{k/m_{\mathrm r}}$ for each bond, and dividing one by the
other cancels $k$:
\begin{equation*}
   \frac{\omega_{\mathrm{OD}}}{\omega_{\mathrm{OH}}}
   = \sqrt{\frac{m_{\mathrm{r},\mathrm{OH}}}{m_{\mathrm{r},\mathrm{OD}}}}
   = \sqrt{\frac{0.9483}{1.7888}} = 0.728 :
\end{equation*}
the O--D vibration has 0.728 times the angular frequency, and its period
is longer by the factor $1/0.728 = 1.37$.

\solutionto{ex:os-coupling}The middle spring is stretched by $x_2 - x_1$,
so $m\ddot{x}_1 = -kx_1 + k'(x_2 - x_1) = -(k + k')x_1 + k'x_2$ and
$m\ddot{x}_2 = -k'(x_2 - x_1) - kx_2 = k'x_1 - (k + k')x_2$. Written as
$m\ddot{\vb{x}} = -\boldsymbol{\Phi}\vb{x}$, the matrix $\boldsymbol{\Phi}$
has rows $(k + k', -k')$ and $(-k', k + k')$. Acting on $(1, 1)$ it gives
$(k, k) = k\,(1, 1)$, so by \cref{eq:os-eigen} $m\omega^2 = k$; acting on $(1, -1)$ it gives $(k + 2k', -k - 2k') = (k +
2k')(1, -1)$, so $m\omega^2 = k + 2k'$.

\solutionto{ex:os-triatomic}Write $\rho = 2m_{\mathrm O}/m_{\mathrm C}$, so
$\vb{u} = (1, -\rho, 1)$. Row by row,
\begin{equation*}
   \boldsymbol{\Phi}\vb{u} = k\,\bigl(1 + \rho,\ -1 - 2\rho - 1,\ \rho + 1\bigr)
   = k\,(1 + \rho)\,(1, -2, 1) ,
\end{equation*}
and, since $\rho\,m_{\mathrm C} = 2m_{\mathrm O}$, $\vb{M}\vb{u} =
(m_{\mathrm O}, -\rho\,m_{\mathrm C}, m_{\mathrm O}) = m_{\mathrm O}(1, -2,
1)$. Hence $\boldsymbol{\Phi}\vb{u} = [k(1 + \rho)/m_{\mathrm O}]\,\vb{M}\vb{u}$,
and $\omega^2 = (k/m_{\mathrm O})(1 + 2m_{\mathrm O}/m_{\mathrm C})$.

\solutionto{ex:os-pair}Dividing each entry by $\sqrt{m_am_b}$, the
mass-weighted Hessian has rows $(k/m_1, -k/\sqrt{m_1m_2})$ and
$(-k/\sqrt{m_1m_2}, k/m_2)$. Its characteristic equation is
\begin{equation*}
   \Bigl(\frac{k}{m_1} - \lambda\Bigr)\Bigl(\frac{k}{m_2} - \lambda\Bigr)
   - \frac{k^2}{m_1m_2} = \lambda^2 - \lambda\,k\Bigl(\frac{1}{m_1}
   + \frac{1}{m_2}\Bigr) = \lambda\Bigl(\lambda - \frac{k}{m_{\mathrm r}}\Bigr) = 0 ,
\end{equation*}
multiplying out, cancelling the two terms $k^2/(m_1m_2)$, taking out the
factor $\lambda$ and using $1/m_{\mathrm r} = 1/m_1 + 1/m_2$. The
eigenvalues are the squared angular frequencies, so $\lambda = 0$ is the
pair moving together without stretching the spring, and $\lambda =
k/m_{\mathrm r}$, $\omega = \sqrt{k/m_{\mathrm r}}$, is the vibration of
\cref{sec:os-bodies}.

\solutionto{ex:os-stiffness}$\omega = 2\pi\times2.99792458\times10^{-5}
\times3657 = \SI{0.6889}{\radian\per\femto\second}$. For a spring
$\omega^2 = k/m_{\mathrm r}$, which, with $k$ in
\si{\electronvolt\per\angstrom\squared}, $m_{\mathrm r}$ in \si{\amu} and
$\omega$ in \si{\radian\per\femto\second}, reads $\omega^2 =
k\times9.648533\times10^{-3}/m_{\mathrm r}$, as for lithium in
\cref{sec:os-small}: the factor of \cref{sec:ne-atoms} turns a force
divided by a mass into an acceleration. Solving for $k$,
\begin{equation*}
   k = \frac{0.9483\times0.6889^2}{9.648533\times10^{-3}}
   = \SI{46.6}{\electronvolt\per\angstrom\squared} ,
\end{equation*}
nearly fifty times stiffer than the lithium atom's hollow of
\cref{sec:os-small}.

\solutionto{ex:os-step}The shortest period in \cref{tab:os-vibrations} for
water is \SI{8.88}{\femto\second}, so the step could be at most
\SI{0.888}{\femto\second}, and a picosecond, \SI{1000}{\femto\second},
would take $1000/0.888 = 1126.1$ steps of that length; 1126 steps fall
just short, so it takes 1127.
